\documentclass[11pt,a4paper]{report}

\usepackage[utf8]{inputenc}
\usepackage[T1]{fontenc}
\usepackage{geometry}
\usepackage{hyperref}
\usepackage{listings}
\usepackage{xcolor}
\usepackage{titlesec}
\usepackage{fancyhdr}
\usepackage{graphicx}
\usepackage{booktabs}
\usepackage{amsmath}

\geometry{margin=2.5cm}

% Code listing style
\lstdefinestyle{uaos}{
    backgroundcolor=\color{black!95},
    basicstyle=\ttfamily\small\color{green!80!black},
    keywordstyle=\color{cyan},
    commentstyle=\color{gray},
    stringstyle=\color{yellow!70!white},
    numbers=left,
    numberstyle=\tiny\color{gray},
    breaklines=true,
    frame=single,
    rulecolor=\color{gray!40},
}
\lstset{style=uaos}

% Header/footer
\pagestyle{fancy}
\fancyhf{}
\lhead{\textsc{Ultimate Amiga OS}}
\rhead{\textsc{Technical Reference Manual}}
\cfoot{\thepage}

\hypersetup{
    colorlinks=true,
    linkcolor=blue,
    urlcolor=cyan,
    pdftitle={UAOS Technical Reference Manual},
    pdfauthor={UAOS Development Team},
}

% -------------------------------------------------------------------------
\begin{document}
% -------------------------------------------------------------------------

\begin{titlepage}
    \centering
    \vspace*{3cm}
    {\Huge\bfseries Ultimate Amiga OS\\[0.5em]}
    {\Large\itshape Technical Reference Manual\\[2em]}
    {\large Version 0.1.0-dev\\[0.5em]}
    {\large \today\\[3em]}
    \rule{\linewidth}{0.5pt}\\[1em]
    {\normalsize
        UAOS is a standalone, bare-metal 64-bit operating system for x86\_64 platforms.\\
        It is derived from the AROS open-source Amiga replacement kernel and features\\
        transparent M68k JIT emulation and native GNU tool integration.
    }
    \vfill
    {\small Copyright \textcopyright\ 2026 UAOS Development Team.
     Released under the AROS Public License (APL).}
\end{titlepage}

\tableofcontents
\listoffigures
\listoftables

% =========================================================================
\chapter{Introduction}
% =========================================================================

\section{Project Overview}

Ultimate Amiga OS (UAOS) is a standalone bare-metal x86\_64 operating system
designed to provide a native Amiga-compatible environment on modern hardware
without requiring any host operating system.  It achieves this through:

\begin{itemize}
    \item An extended AROS microkernel providing Exec, DOS, and Intuition
          semantics on 64-bit hardware.
    \item A transparent M68k JIT emulation sandbox that executes legacy
          Amiga ``Hunk'' binaries inside a 4~GB virtual address space.
    \item An illegal-opcode trap mechanism (\texttt{ILLEGAL}/\texttt{0x4AFC})
          that intercepts system calls at the ROM vector level and dispatches
          them to native 64-bit C implementations.
    \item An MMU sandbox using 2~MB huge pages with hardware-register
          windows deliberately left non-present to transparently virtualise
          custom chipset I/O via \#PF exceptions.
\end{itemize}

\section{Architectural Constraints}

\begin{enumerate}
    \item \textbf{No Linux kernel.}  UAOS boots directly into its own
          x86\_64 AROS-derived microkernel.  No Linux, no POSIX kernel layer.

    \item \textbf{Trap-based thunking.}  M68k Hunk binaries are loaded into
          the 4~GB sandbox.  System calls are intercepted via
          \texttt{ILLEGAL + \$414D} magic-token sequences routed to native
          64-bit AROS libraries.

    \item \textbf{MMU hardware virtualisation.}  x86\_64 page tables map the
          sandbox with 2~MB huge pages.  The Amiga hardware register window
          \texttt{0x00B00000}--\texttt{0x00DFFFFF} is marked \textsc{not present}.

    \item \textbf{Coexisting filesystem namespace.}  \texttt{dos.library}
          handlers resolve Amiga assigns (\texttt{SYS:}, \texttt{S:},
          \texttt{C:}, \texttt{RAM:}) alongside POSIX translation paths.

    \item \textbf{AROS Kickstart firmware.}  The open-source AROS replacement
          ROM is used as the modular firmware base with vectors patched at
          build time.
\end{enumerate}

% =========================================================================
\chapter{System Architecture}
% =========================================================================

\section{Boot Sequence}

\begin{enumerate}
    \item GRUB2 loads the ELF64 kernel via the Multiboot2 protocol.
    \item \texttt{uaos\_kernel\_entry.asm} transitions the CPU from 32-bit
          protected mode into 64-bit long mode, sets up a bootstrap identity-map
          page table, and calls \texttt{uaos\_kernel\_main()}.
    \item The C kernel entry point validates the Multiboot2 magic value,
          initialises the VGA text-mode console and UART, then calls:
          \begin{itemize}
              \item \texttt{APIC\_Init()} --- configures the local APIC for ExtINT delivery.
              \item \texttt{FB\_Init()} --- parses the Multiboot2 framebuffer tag.
              \item \texttt{IDT\_Init()} --- installs 256-vector IDT and remaps 8259A PIC.
              \item \texttt{UAOS\_MMU\_Init()} --- installs the sandbox page tables.
              \item \texttt{RTC\_Init()} --- CMOS real-time clock with UIE interrupt.
              \item \texttt{VFS\_Init()} --- virtual filesystem layer and RAM: mount.
              \item \texttt{BlockDev\_Init()} --- block device layer.
              \item \texttt{virtio\_blk\_init()} --- VirtIO block device detection.
              \item \texttt{ide\_init()} --- IDE/ATAPI controller.
              \item \texttt{UAOS\_ROM\_RegisterAll()} --- populates the ROM module registry.
              \item \texttt{UAOS\_Bridge\_Init()} --- allocates the 4~GB guest RAM window.
              \item \texttt{Task\_Init()} --- Ring-3 task scheduler and TSS.
              \item \texttt{PS2Mouse\_Init()} and \texttt{PS2Kbd\_Init()} --- PS/2 input drivers.
              \item \texttt{net\_stack\_init()} --- TCP/IP stack and NIC auto-probe.
              \item \texttt{Desktop\_Draw()} and \texttt{ShellWin\_Init()} --- graphical shell.
              \item \texttt{ShellWin\_RunStartupSequence()} --- executes \texttt{S:Startup-Sequence}.
          \end{itemize}
    \item Enters the \texttt{hlt}-loop event dispatcher.
\end{enumerate}

\section{M68k Thunk Dispatch}

The thunk mechanism replaces selected Exec jump-table vectors in the AROS ROM
with 8-byte breakout stubs assembled by the \texttt{MK\_THUNK} macro in
\texttt{emulation/rom\_patches/rom\_traps.s}:

\begin{lstlisting}[language={[Motorola68k]Assembler}, caption={MK\_THUNK macro}]
MK_THUNK MACRO
        ILLEGAL             ; $4AFC --- triggers JIT breakout
        dc.w    \1          ; UAOS magic token ($414D = "AM")
        dc.w    \2          ; function index (1-based)
        rts                 ; return to guest after dispatch
        ENDM
\end{lstlisting}

When the JIT core decodes \texttt{0x4AFC} it invokes
\texttt{UAOS\_Bridge\_IllegalOpcode()}.  The C dispatcher:
\begin{enumerate}
    \item Reads the two words immediately following the ILLEGAL opcode.
    \item Validates the \texttt{\$414D} signature.
    \item Extracts arguments from the M68k register state (A1, D0, D1, \ldots).
    \item Calls the native 64-bit stub.
    \item Writes return values back into D0/A0.
    \item Advances the guest PC by 6 bytes.
\end{enumerate}

\section{MMU Sandbox}

\texttt{kernel/exec/mmu\_sandbox.c} constructs four-level x86\_64 paging
tables mapping the entire 4~GB guest address space using 2~MB huge pages.
Table~\ref{tab:memmap} shows the critical address assignments.

\begin{table}[h]
    \centering
    \begin{tabular}{lll}
        \toprule
        \textbf{Address Range} & \textbf{Size} & \textbf{Description} \\
        \midrule
        \texttt{0x00000000--0x001FFFFF} & 2 MB  & Chip RAM \\
        \texttt{0x00200000--0x009FFFFF} & 8 MB  & Fast RAM (lower) \\
        \texttt{0x00A00000--0x00AFFFFF} & 1 MB  & Slow RAM / ranger \\
        \texttt{0x00B00000--0x00DFFFFF} & 3 MB  & \textbf{Custom chip registers (NOT PRESENT)} \\
        \texttt{0x00E00000--0x00EFFFFF} & 1 MB  & Extended ROM / card space \\
        \texttt{0x00F80000--0x00FFFFFF} & 512 KB & Kickstart ROM mirror \\
        \texttt{0x01000000--0x7FFFFFFF} & 2 GB  & Extended Fast RAM \\
        \texttt{0x80000000--0xFFFFFFFF} & 2 GB  & x86\_64 kernel space \\
        \bottomrule
    \end{tabular}
    \caption{UAOS Guest Address Space Map}
    \label{tab:memmap}
\end{table}

Pages in the hardware register window have \texttt{PAGE\_PRESENT} cleared and
\texttt{PAGE\_NO\_CACHE} set.  Any access faults to this region are caught by
the \#PF handler in \texttt{kernel/exec/page\_fault\_handler.c} and forwarded
to the chip emulator.

\section{Ring-3 Userspace and System Calls (\texttt{INT 0x80})}

UAOS implements a Ring-3 (user-mode) task execution model for native 64-bit ELF binaries. These programs are compiled with \texttt{-ffreestanding -nostdlib -fPIE -pie} and linked against \texttt{uaos\_start.o} to produce position-independent binaries wrapped in a 32-byte \texttt{UAOS} header (type \texttt{0x0003} / \texttt{UAOS\_TYPE\_X64}).

\begin{enumerate}
    \item \textbf{GDT \& TSS Setup} --- The ELF64 loader in \texttt{task.c} sets up the task stack, GDT user code/data segments (\texttt{SEG\_USER\_CS}, \texttt{SEG\_USER\_DS}), and registers the task with the Task State Segment (TSS). Transition to Ring 3 is performed via an \texttt{iretq} sequence targeting the binary's entry point.
    \item \textbf{INT 0x80 System Calls} --- Userspace utilities interact with the kernel using software interrupts via the \texttt{INT 0x80} vector. Arguments are passed via registers:
          \begin{itemize}
              \item \textbf{RAX} --- System call number (e.g. \texttt{0x02} for write, \texttt{0x0C} for getcwd).
              \item \textbf{RDI, RSI, RDX} --- Arguments 1, 2, 3.
              \item \textbf{RAX} --- Return value (negative values represent errors).
          \end{itemize}
          The dispatcher in \texttt{syscall\_dispatch.c} implements page allocation (\texttt{sys\_alloc}), process management (\texttt{sys\_spawn}, \texttt{sys\_wait}, \texttt{sys\_exit}), VFS file/directory interfaces (\texttt{sys\_getcwd}, \texttt{sys\_opendir}, \texttt{sys\_readdir}, \texttt{sys\_closedir}, \texttt{sys\_stat}), and GUI windowing syscalls (\texttt{create\_window}, \texttt{destroy\_window}, \texttt{draw\_text}, \texttt{draw\_rect}, \texttt{present}, \texttt{get\_event}, \texttt{set\_scroll\_info}, \texttt{set\_scroll}).
    \item \textbf{Current Userspace Programs} --- \texttt{hello}, \texttt{pwd}, \texttt{file}, \texttt{strings}, \texttt{find}, and \texttt{Guide} are built from \texttt{system/userspace/} and staged into \texttt{C:} (or \texttt{Tools:} for \texttt{Guide}).
    \item \textbf{Execution State \& Redirection} ---
          \begin{itemize}
              \item \textbf{Per-Task CWD} --- Each task struct tracks its own current working directory in \texttt{task\_cwd} which VFS resolves relative paths against.
              \item \textbf{Output Buffering} --- Stdout writes are buffered per-task in \texttt{task\_out} and flushed line-by-line via \texttt{native\_print\_fn} to the active shell window.
              \item \textbf{Synchronous Exit} --- When a child task exits via \texttt{sys\_exit}, the parent task is signaled with \texttt{SIGF\_CHILD}, allowing the command-line interface to wait for command completion.
          \end{itemize}
\end{enumerate}

% =========================================================================
\chapter{Kernel Subsystems}
% =========================================================================

\section{ROM Module Registry (\texttt{rom\_modules.c})}

\texttt{UAOS\_ROM\_RegisterAll()} is called once at boot.  It populates a
static array of \texttt{UaosRomModule} descriptors, each holding:
\begin{itemize}
    \item Library name string (e.g.\ \texttt{"exec.library"}).
    \item Library version number.
    \item 32-bit Amiga base address.
    \item Array of native function pointers (indexed 1-based to match
          \texttt{rom\_traps.s} function indices).
\end{itemize}

\subsection{Registered ROM Modules}

Table~\ref{tab:rommodules} lists the modules registered by
\texttt{UAOS\_ROM\_RegisterAll()}.

\begin{table}[h]
\centering
\begin{tabular}{llcl}
\toprule
\textbf{Module} & \textbf{Version} & \textbf{Base} & \textbf{Purpose} \\
\midrule
\texttt{exec.library} & 45 & 0x00000004 & Process/task control, memory, signals \\
\texttt{utility.library} & 37 & 0x00000050 & String utilities \\
\texttt{console.device} & 40 & 0x00000060 & Console I/O \\
\texttt{mathffp.library} & 40 & 0x00000060 & Soft-float arithmetic \\
\texttt{locale.library} & 38 & 0x000000B0 & Localisation \\
\texttt{ixemul.library} & 53 & 0x00000090 & Unix compatibility layer \\
\texttt{timer.device} & 40 & 0x000000A0 & Timing (RTC-backed) \\
\texttt{keyboard.device} & 40 & 0x000000B0 & Keyboard input (PS/2) \\
\texttt{graphics.library} & 40 & 0x000000C0 & Graphics primitives \\
\texttt{dos.library} & 40 & 0x000000D0 & VFS file system operations \\
\texttt{bsdsocket.library} & 4 & 0x00003000 & BSD socket API \\
\texttt{workbench.library} & 45 & 0x00000000 & Workbench integration \\
\texttt{intuition.library} & 40 & 0x00005000 & GUI/Intuition API \\
\bottomrule
\end{tabular}
\caption{ROM modules registered at boot}
\label{tab:rommodules}
\end{table}

\section{Thunk Handler (\texttt{thunk\_handler.c})}

The \texttt{UAOS\_AMIGA\_TO\_HOST()} macro translates any 32-bit Amiga guest
address to the corresponding host linear address:
\begin{lstlisting}[language=C, caption={Address translation macro}]
#define UAOS_AMIGA_TO_HOST(amiga_addr) \
    ((void *)((uintptr_t)(uaos_ram_base) + (uint32_t)(amiga_addr)))
\end{lstlisting}

\section{Page Fault Handler (\texttt{page\_fault\_handler.c})}

The ISR entry stub \texttt{uaos\_page\_fault\_isr} (written in inline GAS
assembly) saves all 15 GPRs onto the kernel stack, computes pointers to the
\texttt{InterruptFrame} and \texttt{SavedRegs} blocks, and calls the C
function \texttt{UAOS\_PageFaultHandler()}.  On return from the C handler,
GPRs are restored and \texttt{IRETQ} resumes execution past the faulting
instruction.

\section{Native Command System (\texttt{kernel/shell/})}

The shell dispatches \texttt{C:} binaries through a static registry in
\texttt{native\_cmd.c}.  \texttt{k\_native\_cmds[]} maps a lowercase command
name to a \texttt{NativeCmdFn} handler and an optional AmigaDOS-style template
string.  Key components:

\begin{itemize}
    \item \texttt{native\_cmd.c/h} --- registry and case-insensitive lookup.
    \item \texttt{cmd\_template.c/h} --- AmigaDOS-style argument parser
          (\texttt{/A}, \texttt{/K}, \texttt{/S}, \texttt{/N}, \texttt{/M}, \texttt{/F}).
    \item \texttt{resident\_cmd.c/h} --- in-memory resident command cache
          (256~KB, up to 16 slots).
    \item \texttt{cmd\_*.c} --- individual command implementations (65+ commands).
\end{itemize}

\section{Virtual Filesystem Layer (\texttt{vfs.c}, \texttt{ramfs.c})}

UAOS provides a thin VFS dispatch layer over an in-memory RAM filesystem.
Paths use AmigaDOS format: \texttt{VOL:path/to/file}.

\begin{lstlisting}[language=C, caption={VFS file handle}]
typedef struct {
    RamFsNode *node;    /* NULL = invalid / not open */
    uint32_t   pos;     /* current read/write position */
} VfsFile;
\end{lstlisting}

\textbf{RAM Filesystem} (\texttt{ramfs.c}): BSS-backed storage with up to 1024
nodes, 512~KB per file, and a 4~MB data pool.  Auto-mounted at boot as
\texttt{RAM:} with \texttt{T}, \texttt{ENV}, \texttt{CLIPS}, and \texttt{S}
directories.

\textbf{Filesystem drivers}: FAT32 (mount, open, seek, size, format;
cluster read/write pending), PFS3 (mount), EXT4 (read-only mount), and
ISO9660 (CD-ROM reader that loads files into RAMFS at boot).

\textbf{Partition volumes}: The block device layer detects MBR partitions
on VirtIO/IDE devices and registers them as child block devices
(e.g.\ \texttt{virtio01}).  Formatted partitions are auto-mounted in the VFS
by their display name (e.g.\ \texttt{DH0:}) and by their FAT32 volume label
(e.g.\ \texttt{WORK:}).

\textbf{VFS Operations} include \texttt{VFS\_Open}, \texttt{VFS\_Read},
\texttt{VFS\_Write}, \texttt{VFS\_MkDir}, \texttt{VFS\_Delete},
\texttt{VFS\_ResolveDir}, \texttt{VFS\_MountPartition},
\texttt{VFS\_GetMountCount}, and \texttt{VFS\_GetMountName}.

\section{Block Device Layer (\texttt{blockdev.c})}

Unified interface for storage devices.  Key operations:

\begin{itemize}
    \item \texttt{BlockDev\_Register(dev)} / \texttt{BlockDev\_Unregister(dev)}
    \item \texttt{BlockDev\_RegisterPartition(parent, idx, start, count, name)}
          --- creates child partition devices from MBR entries.
    \item \texttt{BlockDev\_CheckFormatted(dev)} --- reads boot sector for
          valid FAT BPB signature.
    \item \texttt{BlockDev\_ReadVolLabel(dev, buf, max)} --- extracts the
          11-byte FAT32 volume label from boot sector offset~71.
    \item \texttt{BlockDev\_Read} / \texttt{BlockDev\_Write} / \texttt{BlockDev\_GetCapacity}
\end{itemize}

\textbf{VirtIO Block Driver} (\texttt{virtio\_blk.c}): PCI scanning for
vendor~0x1AF4 / device~0x1001, capacity reporting from configuration space.
Virtqueue I/O is pending implementation.

% =========================================================================
\chapter{Build System}
% =========================================================================

\section{Prerequisites}

\begin{lstlisting}[language=bash, caption={Install build dependencies (Ubuntu/Debian)}]
sudo apt install nasm gcc binutils xorriso \
    grub-pc-bin grub-common \
    binutils-m68k-linux-gnu python3
\end{lstlisting}

\section{Building the ISO}

\begin{lstlisting}[language=bash, caption={Full build pipeline}]
# Single command builds everything:
./scripts/build_iso.sh

# Clean rebuild:
./scripts/build_iso.sh --clean
\end{lstlisting}

The script compiles the NASM and C kernel sources, links the ELF64 kernel,
embeds any M68k binaries in \texttt{emulation/binaries/}, builds native x86-64
Ring-3 userspace programs from \texttt{system/userspace/}, stages the
\texttt{system/} skeleton into a dynamic \texttt{SYS\_ROOT}, and invokes
\texttt{grub-mkrescue} to produce \texttt{build/Ultimate\_Amiga\_OS.iso}.

\section{Testing in QEMU}

Copy the OVMF variables file once (required for UEFI variable store):

\begin{lstlisting}[language=bash, caption={One-time OVMF setup}]
cp /usr/share/OVMF/OVMF_VARS_4M.fd /tmp/ovmf_vars.fd
\end{lstlisting}

\begin{lstlisting}[language=bash, caption={QEMU test invocation}]
qemu-system-x86_64 \
  -machine q35,usb=off \
  -drive if=pflash,format=raw,readonly=on,\
file=/usr/share/OVMF/OVMF_CODE_4M.fd \
  -drive if=pflash,format=raw,file=/tmp/ovmf_vars.fd \
  -device piix3-ide,id=ide \
  -drive if=none,id=cdrom,media=cdrom,\
file=build/Ultimate_Amiga_OS.iso \
  -device ide-cd,drive=cdrom,bus=ide.0 \
  -m 512M -vga virtio -no-reboot -no-shutdown \
  -serial stdio
\end{lstlisting}

The \texttt{usb=off} flag is required to prevent QEMU's USB tablet device
from conflicting with the PS/2 relative mouse driver. The explicit
\texttt{piix3-ide} controller is needed because the Q35 chipset has no
built-in IDE controller. Add \texttt{-serial stdio} to see kernel debug
output on the terminal.

% =========================================================================
\chapter{Graphics and Window Manager}
% =========================================================================

\section{Framebuffer Initialisation}

GRUB2 fills the Multiboot2 framebuffer tag (type~8) from the UEFI GOP driver.
\texttt{FB\_Init()} parses this tag to populate the global \texttt{FbState g\_fb}
structure, which holds the physical base address, pixel dimensions, pitch
(bytes per row), and bit depth (24 or 32 bpp).

All drawing primitives (\texttt{FB\_FillRect}, \texttt{FB\_DrawHLine},
\texttt{FB\_DrawVLine}, \texttt{FB\_PutPixel}, \texttt{FB\_PutStr}) clip
coordinates to the framebuffer bounds before writing, supporting windows
that extend partially off-screen.

\section{Desktop}

\texttt{Desktop\_Draw()} renders the full Workbench-style backdrop:
\begin{itemize}
    \item \textbf{Menu bar} (top, 20 px) --- Workbench menu labels and clock area.
    \item \textbf{Backdrop} --- light grey base with a dark-grey stipple dot overlay,
          one dot per two pixels alternating rows/columns (classic Amiga pattern).
    \item \textbf{Disk icons} --- top-right corner, stacked vertically.
    \item \textbf{Information window} --- centred boot status window.
    \item \textbf{Status bar} (bottom, 18 px) --- kernel idle message and RAM size.
\end{itemize}

\texttt{Desktop\_RedrawRect(rx, ry, rw, rh)} repaints a sub-rectangle of the
desktop (backdrop + overlays) without a full-screen repaint.  This is used
by the window manager to erase the old window footprint during drag and resize,
avoiding the full-screen flicker that would otherwise occur.

\section{Window Manager}

\texttt{kernel/display/wm.c} implements a lightweight window manager supporting
up to eight windows simultaneously.

\subsection{Data Structures}

\begin{lstlisting}[language=C, caption={WmWindow structure}]
typedef struct {
    int       x, y, w, h;
    char      title[32];
    WM_DrawFn draw;     /* repaint callback: (x, y, w, h) */
    WM_KeyFn  on_key;   /* keystroke callback: (char c)   */
    int       active;
} WmWindow;
\end{lstlisting}

\subsection{Z-Order and Focus}

Windows are stored in \texttt{g\_wins[]} and indexed by a separate
\texttt{g\_zorder[]} array (\texttt{[0]} = back, \texttt{[g\_nwins-1]} = top).
\texttt{WM\_Redraw()} paints windows back-to-front.  Clicking a window calls
\texttt{raise\_window()} to promote it to the top of the z-stack.

\subsection{Drag and Resize}

\begin{itemize}
    \item \textbf{Drag} --- a mouse-down on the title bar records the grab offset.
          Each subsequent mouse move erases the old footprint with
          \texttt{Desktop\_RedrawRect}, updates \texttt{w->x/y}, and repaints
          all windows.
    \item \textbf{Resize} --- a mouse-down on the 16$\times$16 resize grip in the
          bottom-right corner records the base size and origin.  Mouse movement
          computes the new dimensions relative to the drag start point.
    \item Windows may extend off the left, right, and bottom screen edges.
          The title bar is always constrained to remain on-screen.
\end{itemize}

\section{Software Cursor}

\texttt{kernel/display/cursor.c} implements a 16$\times$16 Amiga-style arrow
sprite cursor.  The cursor is rendered by saving the background pixels beneath
it, drawing the sprite, and on each move restoring the saved background before
re-saving and re-drawing at the new position.  This ensures the desktop and
window content are never permanently overwritten by cursor movement.

\section{Additional Windows}

Besides the main shell and Workbench, the kernel provides several standalone
windows:

\begin{itemize}
    \item \textbf{About UAOS} --- \texttt{kernel/display/about\_win.c}
    \item \textbf{Calculator} --- \texttt{kernel/display/calc\_win.c}, opened by \texttt{calculator}
    \item \textbf{Clock} --- \texttt{kernel/display/clock\_win.c}, opened by \texttt{clock}
    \item \textbf{NetInfo} --- \texttt{kernel/display/netinfo\_win.c}, opened by \texttt{netinfo}
    \item \textbf{Pointer preferences} --- \texttt{kernel/display/pointer\_prefs.c}, opened by \texttt{pointer}
    \item \textbf{User window / Guide viewer} --- \texttt{kernel/display/user\_window.c}, opened by \texttt{Guide}
    \item \textbf{VIM editor} --- \texttt{kernel/display/vim\_win.c}, opened by \texttt{vim}
    \item \textbf{File browser} --- \texttt{kernel/display/filebrowser.c}, opened from Workbench icons
\end{itemize}

% =========================================================================
\chapter{Shell and Built-in Commands}
% =========================================================================

The shell window (\texttt{shell\_win.c}) registers with the window manager
via \texttt{WM\_AddWindow()} and implements a scrollable terminal with
up to 128 lines of history and an input line with backspace support.

\section{Command Execution Flow}

The shell resolves a command in the following order:

\begin{enumerate}
    \item \textbf{Shell built-ins} --- implemented directly in the shell.
    \item \textbf{Resident commands} --- cached in memory by \texttt{resident}.
    \item \textbf{Native C: commands} --- dispatched from \texttt{k\_native\_cmds[]}.
    \item \textbf{M68k Hunk binaries} --- loaded into the Musashi emulator.
    \item \textbf{Ring-3 userspace ELF64 binaries} --- run with \texttt{INT 0x80} syscalls.
\end{enumerate}

All command lookup is case-insensitive.  Native commands may declare an
AmigaDOS-style template (\texttt{/A} required, \texttt{/K} keyword,
\texttt{/S} switch, \texttt{/N} numeric, \texttt{/M} multiple, \texttt{/F}
free-form) that is parsed automatically before the handler is invoked.

\section{Built-in Commands}

\begin{table}[h]
    \centering
    \begin{tabular}{ll}
        \toprule
        \textbf{Command} & \textbf{Description} \\
        \midrule
        \texttt{help} & List available commands \\
        \texttt{cd [path]} & Change or show current directory \\
        \texttt{alias [name cmd]} & Create or list aliases \\
        \texttt{unalias <name>} & Remove an alias \\
        \texttt{set [name val]} & Set or list local variables \\
        \texttt{unset <name>} & Remove a local variable \\
        \texttt{path [dirs...]} & Show or set command search path \\
        \texttt{setenv <name> <value>} & Set a global variable \\
        \texttt{unsetenv <name>} & Remove a global variable \\
        \texttt{showconfig} & Show hardware configuration \\
        \bottomrule
    \end{tabular}
    \caption{UAOS Shell Built-in Commands}
    \label{tab:builtins}
\end{table}

\section{Native C: Commands}

\begin{table}[h]
    \centering
    \begin{tabular}{ll}
        \toprule
        \textbf{Command} & \textbf{Description} \\
        \midrule
        \texttt{version} & Show OS version and architecture \\
        \texttt{mem} & Display memory information \\
        \texttt{libs} & List loaded ROM libraries \\
        \texttt{clear} & Clear shell history \\
        \texttt{reboot} & Warm reboot \\
        \texttt{dir [path]} & List directory contents \\
        \texttt{makedir <path>} & Create a directory \\
        \texttt{delete <path>} & Delete a file or empty directory \\
        \texttt{type <file>} & Display file contents \\
        \texttt{copy <src> <dst>} & Copy a file \\
        \texttt{rename <from> <to>} & Rename or move a file \\
        \texttt{echo <text>} & Print text \\
        \texttt{protect <flags> <path>} & Set attributes ($\pm$r, $\pm$h) \\
        \texttt{attr <path>} & Show file attributes \\
        \texttt{info [device]} & Show mounted disks and volumes \\
        \texttt{date} & Show date and time \\
        \texttt{which <cmd>} & Locate a command \\
        \texttt{disks} & List block devices \\
        \texttt{fdisk <device>} & Partition a block device \\
        \texttt{format <dev> [fs]} & Format a partition (FAT32) \\
        \texttt{pointer} & Open pointer preferences \\
        \texttt{run <cmd> [args]} & Run a command in a new CLI \\
        \texttt{assign [name: target]} & Create or list assigns \\
        \texttt{execute <script>} & Run a script file \\
        \texttt{loadwb} & Launch the Workbench desktop \\
        \texttt{calculator} & Open calculator window \\
        \texttt{clock} & Open clock window \\
        \texttt{ifconfig} & Show or configure network \\
        \texttt{ping} & Send ICMP echo requests \\
        \texttt{route} & Show routing table and ARP cache \\
        \texttt{nslookup} & Resolve a hostname \\
        \texttt{ntpd} & Synchronise time via NTP \\
        \texttt{netstart / netstop} & Start or stop the network stack \\
        \texttt{netinfo} & Open network info window \\
        \texttt{grep / more / list / search / sort / join} & Text/file processing \\
        \texttt{vim} & Inline text editor \\
        \texttt{newcli / newshell} & Open a new shell window \\
        \texttt{ask / prompt / why / failat / quit / endcli} & Shell scripting \\
        \texttt{resident} & Manage resident commands \\
        \texttt{ps / jobs / wait / changetaskpri} & Process/task control \\
        \texttt{filenote / relabel / install / diskchange / addbuffers} & Disk maintenance \\
        \texttt{requestchoice / requestfile} & Dialogs \\
        \texttt{status / avail / stack} & System utilities \\
        \texttt{getenv / unset} & Environment access \\
        \bottomrule
    \end{tabular}
    \caption{UAOS Native C: Commands}
    \label{tab:nativecmds}
\end{table}

\section{Ring-3 Userspace Commands}

\begin{table}[h]
    \centering
    \begin{tabular}{ll}
        \toprule
        \textbf{Command} & \textbf{Description} \\
        \midrule
        \texttt{pwd} & Print working directory \\
        \texttt{find [path] [-name pat] [-type f|d]} & Recursively search directories \\
        \texttt{file <path>...} & Identify file format from magic numbers \\
        \texttt{strings <path>... [-n minlen]} & Extract printable strings \\
        \texttt{hello} & Print a simple greeting \\
        \texttt{Guide} & AmigaGuide help viewer \\
        \bottomrule
    \end{tabular}
    \caption{UAOS Ring-3 Userspace Commands}
    \label{tab:userspacecmds}
\end{table}

\section{Filesystem Operations}

The shell integrates with the VFS layer for all filesystem operations.
Partition volumes (e.g.\ \texttt{DH0:}, \texttt{WORK:}) can be accessed
with \texttt{cd} and \texttt{dir} once auto-mounted at boot or after
formatting with the \texttt{format} command.

% =========================================================================
% =========================================================================
\chapter{TCP/IP Networking}
\label{ch:networking}
% =========================================================================

UAOS includes a complete freestanding TCP/IP stack.  It auto-detects the
network card at boot, attempts DHCP, and falls back to a static IP address
if DHCP times out.  The stack runs entirely in the kernel event loop --- no
threads or blocking system calls are required.

% -------------------------------------------------------------------------
\section{Architecture}
% -------------------------------------------------------------------------

\begin{verbatim}
Shell / bsdsocket.library
        |
        v
 icmp / tcp / udp   (kernel/net/)
        |
        v
  ip.c  <-->  arp.c
        |
        v
 NetDevice  (abstract hardware interface)
        |
   +----+----+
   |         |
   v         v
virtio_net  e1000
(QEMU)   (VirtualBox)
\end{verbatim}

% -------------------------------------------------------------------------
\section{Network Hardware Drivers}
% -------------------------------------------------------------------------

\subsection{NetDevice Abstraction (\texttt{net\_device.h})}

All stack components call through the \texttt{NetDevice} interface rather
than any specific driver.  \texttt{netdev\_probe()} (called by
\texttt{net\_stack\_init\_ex()}) scans PCI for a supported device, registers
the first one found, and calls its \texttt{init()} function.  The e1000 is
tried before virtio-net so that VirtualBox guests are served without
configuration.

\begin{lstlisting}[language=C]
typedef struct NetDevice {
    const char *name;
    int  (*init)(struct NetDevice *dev);
    void (*get_mac)(struct NetDevice *dev, uint8_t *buf);
    int  (*send)(struct NetDevice *dev, const uint8_t *data, uint16_t len);
    void (*poll)(struct NetDevice *dev);
    void (*set_rx_callback)(struct NetDevice *dev, netdev_rx_fn cb);
    void (*setup_irq)(struct NetDevice *dev);
    void *priv;
} NetDevice;
\end{lstlisting}

\subsection{VirtIO-Net (\texttt{virtio\_net.c})}

\begin{itemize}
  \item PCI vendor \texttt{0x1AF4}, device \texttt{0x1000} (legacy VirtIO~0.9)
        or \texttt{0x1041} (modern).
  \item I/O-port or MMIO BAR0; split virtqueue descriptor rings.
  \item MSI-X capability is present but kept disabled; interrupts use the
        8259A~PIC (IRQ~10 on Q35 QEMU).
  \item Used by: QEMU \texttt{-device virtio-net-pci}.
\end{itemize}

\subsection{Intel 82540EM e1000 (\texttt{e1000.c})}

\begin{itemize}
  \item PCI vendor \texttt{0x8086}, device \texttt{0x100E} (and related 8254x family IDs).
  \item 128~KB MMIO BAR0; legacy 16-byte TX/RX descriptor rings (32 slots each).
  \item MAC address read from EEPROM via the EERD register at initialisation.
  \item IRQ delivered via ICR (self-clearing); routed through the 8259A PIC.
  \item Used by: VirtualBox (Intel PRO/1000 MT Desktop) and QEMU
        \texttt{-device e1000}.
\end{itemize}

% -------------------------------------------------------------------------
\section{Stack Initialisation}
% -------------------------------------------------------------------------

Call \texttt{net\_stack\_init()} once from \texttt{uaos\_kernel\_main()}:

\begin{lstlisting}[language=C]
/* DHCP with 1-second timeout, static fallback */
net_stack_init(IPV4(10,0,2,15), IPV4(10,0,2,2), IPV4(255,255,255,0));

/* Explicit DHCP timeout (ms) */
net_stack_init_ex(IPV4(10,0,2,15), IPV4(10,0,2,2),
                  IPV4(255,255,255,0), 2000);
\end{lstlisting}

\noindent The \texttt{IPV4(a,b,c,d)} macro builds a host-byte-order \texttt{ipv4\_t}:

\begin{lstlisting}[language=C]
#define IPV4(a,b,c,d) ((ipv4_t)(  ((uint32_t)(a)<<24) \
                                 | ((uint32_t)(b)<<16) \
                                 | ((uint32_t)(c)<< 8) \
                                 |  (uint32_t)(d)     ))
\end{lstlisting}

% -------------------------------------------------------------------------
\section{Stack API Reference}
% -------------------------------------------------------------------------

\begin{table}[h]
\centering
\begin{tabular}{ll}
\toprule
\textbf{Function} & \textbf{Description} \\
\midrule
\texttt{net\_stack\_init(ip, gw, nm)} & Init; DHCP then static fallback \\
\texttt{net\_stack\_init\_ex(ip, gw, nm, ms)} & Init with explicit DHCP timeout \\
\texttt{net\_stack\_is\_up()} & 1 if NIC found and stack is live \\
\texttt{net\_stack\_dhcp\_used()} & 1 if IP came from DHCP \\
\texttt{net\_stack\_poll()} & Process pending RX frames \\
\texttt{net\_stack\_tick()} & TCP retransmit / timeout tick \\
\texttt{net\_stack\_get\_ip()} & Return assigned IPv4 address \\
\texttt{net\_ip\_to\_str(ip, buf)} & Format \texttt{ipv4\_t} as \texttt{"a.b.c.d"} \\
\texttt{net\_str\_to\_ip(str, out)} & Parse \texttt{"a.b.c.d"} string \\
\bottomrule
\end{tabular}
\caption{TCP/IP stack top-level API (\texttt{kernel/net/stack.h})}
\end{table}

% -------------------------------------------------------------------------
\section{ICMP (Ping)}
% -------------------------------------------------------------------------

\begin{lstlisting}[language=C]
icmp_clear_reply();
icmp_ping(IPV4(10,0,2,2), 1);          /* send echo request seq=1 */
/* poll until reply or timeout */
for (int i = 0; i < 10 && !icmp_got_reply(); i++) {
    net_stack_poll();
    /* sleep ~100 ms */
}
if (icmp_got_reply()) { /* success */ }
\end{lstlisting}

% -------------------------------------------------------------------------
\section{UDP}
% -------------------------------------------------------------------------

\begin{lstlisting}[language=C]
int s = udp_open(0);          /* ephemeral local port */
uint8_t msg[] = "hello";
udp_send(s, IPV4(10,0,2,2), 9, msg, sizeof(msg));

uint8_t buf[512]; ipv4_t src; uint16_t sport;
int n = udp_recv(s, buf, sizeof(buf), &src, &sport);
if (n > 0) { /* buf[0..n-1] contains datagram payload */ }
udp_close(s);
\end{lstlisting}

Up to 8 simultaneous UDP sockets.  Each socket has a 2~KB receive ring
buffer.  \texttt{udp\_recv()} is non-blocking and returns 0 if no datagram
is waiting.

% -------------------------------------------------------------------------
\section{TCP}
% -------------------------------------------------------------------------

\subsection{Client Connection}

\begin{lstlisting}[language=C]
int sock = tcp_connect(IPV4(93,184,216,34), 80, 0);
while (tcp_state(sock) != TCP_ESTABLISHED) net_stack_poll();

uint8_t req[] = "GET / HTTP/1.0\r\nHost: example.com\r\n\r\n";
tcp_send(sock, req, sizeof(req) - 1);

uint8_t buf[512]; int n;
while ((n = tcp_recv(sock, buf, sizeof(buf))) > 0) {
    /* process buf[0..n-1] */
    net_stack_poll();
}
tcp_close(sock);
\end{lstlisting}

\subsection{Passive Server}

\begin{lstlisting}[language=C]
int lsock = tcp_listen(8080);
while (1) {
    net_stack_poll();
    int csock = tcp_accept(lsock);
    if (csock < 0) continue;
    /* handle connection ... */
    tcp_close(csock);
}
\end{lstlisting}

\subsection{TCP API Reference}

\begin{table}[h]
\centering
\begin{tabular}{ll}
\toprule
\textbf{Function} & \textbf{Description} \\
\midrule
\texttt{tcp\_connect(ip, port, lport)} & Active open; socket index or $-1$ \\
\texttt{tcp\_listen(port)} & Passive open; socket index or $-1$ \\
\texttt{tcp\_accept(lsock)} & Accept pending connection \\
\texttt{tcp\_send(sock, data, len)} & Queue data; returns bytes queued \\
\texttt{tcp\_recv(sock, buf, max)} & Non-blocking receive \\
\texttt{tcp\_close(sock)} & Graceful close (sends FIN) \\
\texttt{tcp\_state(sock)} & Returns \texttt{TcpState} enum \\
\texttt{tcp\_tick()} & Retransmit / timeout (call from PIT) \\
\bottomrule
\end{tabular}
\caption{TCP API (\texttt{kernel/net/tcp.h})}
\end{table}

\noindent Up to 8 simultaneous TCP connections; 4~KB TX and 4~KB RX ring
buffers per socket.

% -------------------------------------------------------------------------
\section{DHCP}
% -------------------------------------------------------------------------

\texttt{net\_stack\_init\_ex()} invokes DHCP automatically.  For direct
use:

\begin{lstlisting}[language=C]
DhcpLease lease;
if (dhcp_request(&lease, 2000)) {
    /* lease.ip, .gateway, .netmask, .dns filled */
}
\end{lstlisting}

% -------------------------------------------------------------------------
\section{Polling Model}
% -------------------------------------------------------------------------

The stack has no dedicated network thread.  RX is processed via two paths:

\begin{enumerate}
  \item \textbf{IRQ path} --- the NIC fires an interrupt; the handler calls
        \texttt{e1000\_poll()} / \texttt{virtio\_net\_poll()}, which invokes
        the RX callback, which drives \texttt{eth\_rx()} $\to$ \texttt{arp\_rx()} /
        \texttt{ip\_rx()} $\to$ TCP/UDP demux.
  \item \textbf{Poll path} --- \texttt{net\_stack\_poll()} is called from the
        main \texttt{hlt} loop and from \texttt{CMD\_YIELD()} inside shell commands.
\end{enumerate}

Both paths call \texttt{netdev\_poll()} which is guarded by a per-driver
reentrancy lock, making simultaneous calls safe.

% -------------------------------------------------------------------------
\section{Shell Network Commands}
% -------------------------------------------------------------------------

\begin{table}[h]
\centering
\begin{tabular}{ll}
\toprule
\textbf{Command} & \textbf{Description} \\
\midrule
\texttt{netstart} & Initialise stack from \texttt{S:net.conf} \\
\texttt{netstop} & Shut down stack, release DHCP lease \\
\texttt{ifconfig [dhcp | <ip> <gw>]} & Show or configure network \\
\texttt{ping <host> [count]} & Send ICMP echo requests \\
\texttt{route} & Display routing table and ARP cache \\
\texttt{nslookup <host> [server]} & Resolve hostname via DNS \\
\texttt{ntpd [server]} & Synchronise time via NTP \\
\texttt{netinfo} & Open network info window \\
\bottomrule
\end{tabular}
\caption{Network shell commands}
\end{table}

\subsection{Network Configuration Files}

\begin{table}[h]
\centering
\begin{tabular}{lll}
\toprule
\textbf{File} & \textbf{Purpose} & \textbf{Default} \\
\midrule
\texttt{S:net.conf} & Mode (dhcp/static), IP, netmask, gateway, DNS & dhcp / 10.0.2.15/24 \\
\texttt{S:ntp.conf} & NTP server hostname & pool.ntp.org \\
\texttt{S:timezone.conf} & IANA timezone name & Australia/Sydney \\
\bottomrule
\end{tabular}
\caption{Network configuration files}
\end{table}

\subsection{\texttt{bsdsocket.library} LVO Table}

The \texttt{bsdsocket.library} ROM module exposes a standard AmigaOS BSD
socket LVO table to M68k binaries.  Descriptors 0--7 map to TCP sockets and
8--15 to UDP sockets.

\begin{table}[h]
\centering
\begin{tabular}{cll}
\toprule
\textbf{LVO} & \textbf{Function} & \textbf{Arguments} \\
\midrule
-30 & \texttt{socket} & domain/D0, type/D1, protocol/D2 \\
-36 & \texttt{bind} & fd/D0, sockaddr/A0, addrlen/D1 \\
-42 & \texttt{listen} & fd/D0, backlog/D1 \\
-48 & \texttt{accept} & fd/D0, sockaddr/A0, addrlen/A1 \\
-54 & \texttt{connect} & fd/D0, sockaddr/A0, addrlen/D1 \\
-60 & \texttt{send} & fd/D0, buf/A0, len/D1, flags/D2 \\
-66 & \texttt{sendto} & fd/D0, buf/A0, len/D1, flags/D2, addr/A1, addrlen/D3 \\
-72 & \texttt{recv} & fd/D0, buf/A0, len/D1, flags/D2 \\
-78 & \texttt{recvfrom} & fd/D0, buf/A0, len/D1, flags/D2, addr/A1, addrlen/A2 \\
-84 & \texttt{closesocket} & fd/D0 \\
-96 & \texttt{setsockopt} & fd/D0, level/D1, optname/D2, val/A0, optlen/D3 \\
-102 & \texttt{getsockopt} & fd/D0, level/D1, optname/D2, val/A0, optlen/A1 \\
-108 & \texttt{IoctlSocket} & fd/D0, req/D1, arg/A0 \\
-132 & \texttt{inet\_addr} & str/A0 \\
-138 & \texttt{inet\_ntoa} & addr/D0 \\
-210 & \texttt{gethostbyname} & name/A0 \\
\bottomrule
\end{tabular}
\caption{\texttt{bsdsocket.library} LVO table}
\end{table}

% =========================================================================
\chapter{Troubleshooting}
% =========================================================================

\section{GRUB: ``Invalid arch-dependent ELF magic''}

This error means GRUB received a file that is not a valid ELF64 executable.
Ensure \texttt{uaos-kernel.elf} was produced by the full NASM + GCC + LD
pipeline (not the Python mock generator used for structural validation only).

\section{GRUB: ``You need to load the kernel first''}

GRUB could not find the \texttt{multiboot2} header within the first 32~KB of
the binary.  Verify the linker script places \texttt{.multiboot2} first and
that the NASM source emits the section correctly.

\section{Kernel halts immediately after boot}

Enable serial output and check the UART log:
\begin{lstlisting}[language=bash]
qemu-system-x86_64 -cdrom build/Ultimate_Amiga_OS.iso \
    -m 512M -boot d -serial stdio 2>&1 | tee boot.log
\end{lstlisting}

% =========================================================================
\appendix
\chapter{Amiga Memory Map Reference}
% =========================================================================

See Table~\ref{tab:memmap} in Chapter~2.  The full OCS/ECS/AGA custom chip
register map follows Commodore's \emph{Hardware Reference Manual, 3rd Ed.}

\end{document}
